% ============================================================================
% STRUCTURE COMPLÈTE DU RAPPORT PFE INGÉNIEUR - 6 MOIS
% Projet: Plateforme intelligente de configuration et recommandation de 
% produits d'assurance basée sur l'intelligence artificielle
% Entreprise: Vermeg Tunisie
% École: Esprit École d'Ingénieurs
% Méthodologie: Agile/Scrum
% ============================================================================

\documentclass[12pt,a4paper]{report}
\usepackage[utf8]{inputenc}
\usepackage[T1]{fontenc}
\usepackage[french]{babel}
\usepackage{graphicx}
\usepackage{geometry}
\usepackage{titlesec}
\usepackage{tocloft}
\usepackage{hyperref}
\usepackage{listings}
\usepackage{xcolor}
\usepackage{booktabs}
\usepackage{multirow}
\usepackage{caption}

% Configuration des marges
\geometry{
    left=2.5cm,
    right=2.5cm,
    top=2.5cm,
    bottom=2.5cm
}

% ============================================================================
% PAGE DE TITRE
% ============================================================================
\begin{document}

\begin{titlepage}
    \centering
    \vspace*{2cm}
    
    % Logo Esprit
    \includegraphics[width=4cm]{logo_esprit.png}
    
    \vspace{1cm}
    
    % Logo Vermeg
    \includegraphics[width=4cm]{logo_vermeg.png}
    
    \vspace{2cm}
    
    \textbf{\Large Projet de Fin d'Études d'Ingénieur}
    
    \vspace{1cm}
    
    \textbf{\Large Pour l'obtention du diplôme d'Ingénieur en Informatique}
    
    \vspace{1.5cm}
    
    \Huge \textbf{Plateforme Intelligente de Configuration et Recommandation de Produits d'Assurance Basée sur l'Intelligence Artificielle}
    
    \vspace{2cm}
    
    \Large \textbf{Réalisé par:} \\
    [Votre Nom]
    
    \vspace{1cm}
    
    \Large \textbf{Encadré par:} \\
    [Nom Encadreur Académique] - Esprit \\
    [Nom Encadreur Entreprise] - Vermeg Tunisie
    
    \vspace{1.5cm}
    
    \Large \textbf{Soutenu le:} \\
    [Date de soutenance]
    
    \vspace{1cm}
    
    \Large \textbf{Devant le jury composé de:} \\
    Président: [Nom] \\
    Examinateur: [Nom] \\
    Encadreur: [Nom]
    
    \vspace{1cm}
    
    \textbf{Année Universitaire:} 2025/2026
    
\end{titlepage}

% ============================================================================
% REMERCIEMENTS
% ============================================================================
\chapter*{Remerciements}
\addcontentsline{toc}{chapter}{Remerciements}

Je tiens à exprimer ma profonde gratitude à toutes les personnes qui ont contribué à l'aboutissement de ce projet de fin d'études...

% ============================================================================
% DÉDICACES
% ============================================================================
\chapter*{Dédicaces}
\addcontentsline{toc}{chapter}{Dédicaces}

Je dédie ce travail à...

% ============================================================================
% RÉSUMÉ
% ============================================================================
\chapter*{Résumé}
\addcontentsline{toc}{chapter}{Résumé}

Ce projet de fin d'études s'inscrit dans le cadre du développement d'une plateforme intelligente de configuration et recommandation de produits d'assurance basée sur l'intelligence artificielle pour l'entreprise Vermeg Tunisie...

% ============================================================================
% ABSTRACT
% ============================================================================
\chapter*{Abstract}
\addcontentsline{toc}{chapter}{Abstract}

This final year project focuses on the development of an intelligent platform for configuration and recommendation of insurance products based on artificial intelligence for Vermeg Tunisia...

% ============================================================================
% TABLE DES MATIÈRES
% ============================================================================
\tableofcontents

% ============================================================================
% LISTE DES FIGURES
% ============================================================================
\listoffigures

% ============================================================================
% LISTE DES TABLEAUX
% ============================================================================
\listoftables

% ============================================================================
% LISTE DES ABREVIATIONS
% ============================================================================
\chapter*{Liste des Abréviations}
\addcontentsline{toc}{chapter}{Liste des Abréviations}

\begin{tabular}{ll}
    IA & Intelligence Artificielle \\
    API & Application Programming Interface \\
    CI/CD & Continuous Integration/Continuous Deployment \\
    DTO & Data Transfer Object \\
    RAG & Retrieval-Augmented Generation \\
    SaaS & Software as a Service \\
    UI/UX & User Interface/User Experience \\
    \end{tabular}

% ============================================================================
% INTRODUCTION GÉNÉRALE
% ============================================================================
\chapter{Introduction Générale}
\label{chap:introduction}

\section{Contexte du projet}
L'industrie de l'assurance fait face à une transformation numérique majeure...

\section{Problématique}
Comment développer une plateforme intelligente capable de configurer et recommander des produits d'assurance adaptés aux besoins spécifiques des clients tout en assurant une expérience utilisateur optimale?

\section{Objectifs du projet}
\subsection{Objectif général}
Développer une plateforme intelligente de configuration et recommandation de produits d'assurance...

\subsection{Objectifs spécifiques}
\begin{itemize}
    \item Implémenter une architecture microservices moderne
    \item Intégrer un moteur de recommandation basé sur l'IA
    \item Assurer une gestion métier cohérente des produits, packs et garanties
    \item Garantir une sécurité robuste via Keycloak
    \item Déployer une solution DevOps complète (Docker, CI/CD, Monitoring)
\end{itemize}

\section{Méthodologie adoptée}
Ce projet est développé selon la méthodologie Agile/Scrum pour une approche itérative et adaptative...

\section{Plan du rapport}
Ce rapport est structuré comme suit...

% ============================================================================
% CHAPITRE 1: CONTEXTE GÉNÉRAL ET ÉTUDE DE L'EXISTANT
% ============================================================================
\chapter{Contexte Général et Étude de l'Existant}
\label{chap:contexte}

\section{Présentation de l'entreprise d'accueil}
\subsection{Vermeg Tunisie}
Historique, mission, organisation...

\subsection{Le département d'accueil}
Description de l'équipe et du contexte du projet...

\section{Contexte du projet}
\subsection{Problématique métier}
Les défis actuels de l'industrie de l'assurance...

\subsection{Besoins identifiés}
Analyse des besoins fonctionnels et techniques...

\section{État de l'art}
\subsection{Solutions de recommandation IA dans l'assurance}
Analyse des solutions existantes sur le marché...

\subsection{Architectures microservices dans le secteur financier}
Étude des architectures modernes...

\subsection{Technologies IA pour la recommandation}
Overview des technologies disponibles (LLMs, RAG, Vector Search)...

\section{Critique de l'existant}
\subsection{Limites des solutions actuelles}
Analyse des gaps et limitations...

\subsection{Opportunités d'amélioration}
Axes d'innovation proposés...

\section{Conclusion du chapitre}
Synthèse et transition vers le chapitre suivant...

% ============================================================================
% CHAPITRE 2: SPECIFICATION ET ANALYSE DES BESOINS
% ============================================================================
\chapter{Spécification et Analyse des Besoins}
\label{chap:besoins}

\section{Acteurs du système}
\subsection{Acteurs principaux}
Administrateur, Assuré, Agent d'assurance...

\subsection{Acteurs secondaires}
Systèmes externes, API tierces...

\section{Recueil des besoins}
\subsection{Besoins fonctionnels}
\subsubsection{Gestion des produits d'assurance}
CRUD produits, gestion des métadonnées...

\subsubsection{Gestion des packs et garanties}
Configuration des packs, association des garanties...

\subsubsection{Moteur de recommandation IA}
Analyse de profil, scoring, Top 3 recommandations...

\subsubsection{Chatbot intelligent}
Configuration conversationnelle, recommandation interactive...

\subsubsection{Gestion des utilisateurs et authentification}
Keycloak, rôles, permissions...

\subsection{Besoins non fonctionnels}
\subsubsection{Performance}
Temps de réponse, scalabilité...

\subsubsection{Sécurité}
Authentification, autorisation, protection des données...

\subsubsection{Maintenabilité}
Code propre, architecture modulaire, documentation...

\subsubsection{Disponibilité}
Uptime, redondance...

\section{Diagrammes de cas d'utilisation}
\subsection{Diagramme global}
Figure: Cas d'utilisation principaux...

\subsection{Description détaillée des cas d'utilisation}
\subsubsection{CU-01: Créer un produit d'assurance}
Description, préconditions, scénario nominal...

\subsubsection{CU-02: Configurer un pack}
...

\subsubsection{CU-03: Obtenir une recommandation IA}
...

\section{Contraintes du projet}
\subsection{Contraintes techniques}
Stack imposée, compatibilité...

\subsection{Contraintes temporelles}
Calendrier de 6 mois, livrables intermédiaires...

\subsection{Contraintes budgétaires}
Ressources disponibles...

\section{Conclusion du chapitre}
...

% ============================================================================
% CHAPITRE 3: CONCEPTION
% ============================================================================
\chapter{Conception}
\label{chap:conception}

\section{Architecture globale du système}
\subsection{Architecture microservices}
Description de l'approche microservices...

\subsection{Diagramme d'architecture}
Figure: Architecture globale avec API Gateway, Eureka, Services...

\section{Conception de la base de données}
\subsection{Modèle de données MongoDB}
\subsubsection{Entité Produit}
Structure, champs, relations...

\subsubsection{Entité Pack}
Structure, champs, relations...

\subsubsection{Entité Garantie}
Structure, champs, relations...

\subsubsection{Entité PackGarantie}
Champs métier: tauxRemboursement, plafond, delaiCarence, condition, actif, ordreAffichage...

\subsubsection{Diagramme de classes}
Figure: Hiérarchie Produit -> Pack -> Garantie...

\section{Conception des microservices}
\subsection{Service GestionProduit}
Responsabilités, endpoints, architecture interne...

\subsection{Service GestionUser}
...

\subsection{API Gateway}
Routage, authentification, load balancing...

\subsection{Service Discovery (Eureka)}
...

\section{Conception du moteur de recommandation IA}
\subsection{Architecture de recommandation}
Composants: Profil Client, Scoring Engine, LLM Integration...

\subsection{Algorithme de scoring}
Formule détaillée:
\begin{equation}
score = (ageWeight \times 0.15) + (budgetWeight \times 0.20) + (beneficiaryWeight \times 0.15) + (medicalRiskWeight \times 0.20) + (guaranteeMatchWeight \times 0.20) + (coverageLevelWeight \times 0.10)
\end{equation}

\subsection{Intégration Google Gemini API}
Prompt engineering, extraction structurée...

\section{Conception du chatbot}
\subsection{Architecture conversationnelle}
Gestion de la mémoire, contexte, RAG...

\subsection{Flux de dialogue}
Figure: Diagramme de séquence conversationnel...

\section{Conception de la sécurité}
\subsection{Architecture Keycloak}
Realm, clients, rôles, flows d'authentification...

\subsection{Sécurité des microservices}
JWT, OAuth2, Spring Security...

\section{Conception DevOps}
\subsection{Dockerisation}
Structure Dockerfile pour chaque service...

\subsection{CI/CD Pipeline}
GitHub Actions workflow...

\subsection{Monitoring}
Prometheus + Grafana dashboards...

\section{Conclusion du chapitre}
...

% ============================================================================
% CHAPITRE 5: MÉTHODOLOGIE AGILE/SCRUM
% ============================================================================
\chapter{Méthodologie Agile/Scrum}
\label{chap:scrum}

\section{Choix de la méthodologie Agile/Scrum}
\subsection{Justification du choix}
Pourquoi Agile pour ce projet PFE...

\subsection{Adaptation au contexte PFE}
Spécificités de l'approche Agile en milieu académique...

\subsection{Avantages pour ce projet}
Itérations, validation progressive, flexibilité...

\section{Organisation Scrum}
\subsection{Rôles et responsabilités}
Product Owner (encadreur), Scrum Master, Development Team...

\subsection{Artefacts Scrum}
Product Backlog, Sprint Backlog, Increment...

\subsection{Cérémonies Scrum}
Sprint Planning, Daily Scrum, Sprint Review, Sprint Retrospective...

\section{Vue d'ensemble des sprints}
Calendrier des 10 sprints de 2 semaines sur 6 mois...

\section{Suivi de projet}
\subsection{Burndown charts}
Figure: Évolution du travail par sprint...

\subsection{Velocity tracking}
Métriques de vélocité de l'équipe...

\subsection{Gestion des risques}
Identification et mitigation des risques...

\section{Conclusion du chapitre}
...

% ============================================================================
% CHAPITRE 6: SPRINTS 1-2 - INITIALISATION ET ARCHITECTURE
% ============================================================================
\chapter{Sprints 1-2: Initialisation et Architecture Microservices}
\label{chap:sprints1-2}

\section{Vue d'ensemble de la phase}
Objectifs globaux des 4 premières semaines...

\section{Sprint 1: Initialisation et Setup (Semaines 1-2)}
\subsection{Objectifs du sprint}
Configuration environnement, setup projet, architecture de base...

\subsection{Planification}
\subsubsection{User Stories}
US-1.1: Configurer l'environnement de développement \\
US-1.2: Initialiser le projet Spring Boot \\
US-1.3: Initialiser le projet Angular \\
US-1.4: Configurer MongoDB Docker \\
US-1.5: Configurer Keycloak Docker

\subsubsection{Estimation et priorisation}
Story points, priorité des US...

\subsection{Réalisation}
\subsubsection{Configuration de l'environnement}
JDK 21, Node.js, Docker, MongoDB...

\subsubsection{Setup Backend Spring Boot}
Structure Maven, dépendances, configuration de base...

\subsubsection{Setup Frontend Angular}
Structure Angular 18 standalone, PrimeNG, Bootstrap...

\subsubsection{Configuration Docker}
Docker Compose pour MongoDB et Keycloak...

\subsection{Résultats et livrables}
Environnement de développement opérationnel, architecture de base...

\subsection{Difficultés rencontrées et solutions}
Problèmes techniques et résolutions...

\section{Sprint 2: Architecture Microservices et Gateway (Semaines 3-4)}
\subsection{Objectifs du sprint}
Implémenter API Gateway, Service Discovery, architecture de base...

\subsection{Planification}
\subsubsection{User Stories}
US-2.1: Implémenter Eureka Server \\
US-2.2: Configurer Spring Cloud Gateway \\
US-2.3: Définir les routes inter-services \\
US-2.4: Configurer le load balancing

\subsubsection{Estimation et priorisation}
...

\subsection{Réalisation}
\subsubsection{Implémentation Eureka Server}
Configuration du service registry...

\subsubsection{Configuration Spring Cloud Gateway}
Routes, filtres, load balancing...

\subsubsection{Architecture microservices}
Communication inter-services, résilience...

\subsection{Résultats et livrables}
Infrastructure microservices fonctionnelle...

\subsection{Difficultés rencontrées et solutions}
...

\section{Synthèse de la phase}
Bilan des sprints 1-2, leçons apprises...

\section{Conclusion du chapitre}
...

% ============================================================================
% CHAPITRE 7: SPRINTS 3-4 - SÉCURITÉ ET DONNÉES
% ============================================================================
\chapter{Sprints 3-4: Sécurité et Modèle de Données}
\label{chap:sprints3-4}

\section{Vue d'ensemble de la phase}
Objectifs globaux des semaines 5-8...

\section{Sprint 3: Sécurité et Authentification (Semaines 5-6)}
\subsection{Objectifs du sprint}
Intégration Keycloak, Spring Security, JWT...

\subsection{Planification}
\subsubsection{User Stories}
US-3.1: Configurer Keycloak realm et clients \\
US-3.2: Intégrer Spring Security dans les microservices \\
US-3.3: Implémenter l'authentification JWT \\
US-3.4: Configurer les rôles et permissions \\
US-3.5: Implémenter le logout Keycloak

\subsubsection{Estimation et priorisation}
...

\subsection{Réalisation}
\subsubsection{Configuration Keycloak}
Realm, clients, rôles, utilisateurs de test...

\subsubsection{Intégration Spring Security}
Configuration JWT, filtres de sécurité...

\subsubsection{Gestion des rôles et permissions}
Autorisations par endpoint...

\subsection{Résultats et livrables}
Système d'authentification sécurisé...

\subsection{Difficultés rencontrées et solutions}
...

\section{Sprint 4: Modèle de Données MongoDB (Semaines 7-8)}
\subsection{Objectifs du sprint}
Définir et implémenter le modèle de données complet...

\subsection{Planification}
\subsubsection{User Stories}
US-4.1: Définir l'entité Produit \\
US-4.2: Définir l'entité Pack \\
US-4.3: Définir l'entité Garantie \\
US-4.4: Définir l'entité PackGarantie avec champs métier \\
US-4.5: Configurer les repositories MongoDB \\
US-4.6: Créer les DTOs

\subsubsection{Estimation et priorisation}
...

\subsection{Réalisation}
\subsubsection{Entité Produit}
Champs, annotations MongoDB, validation...

\subsubsection{Entité Pack}
Structure, relation avec Produit...

\subsubsection{Entité Garantie}
Définition des garanties...

\subsubsection{Entité PackGarantie}
Champs métier: tauxRemboursement, plafond, delaiCarence, condition, actif, ordreAffichage...

\subsubsection{Configuration MongoDB}
Connection, repositories, indexes...

\subsection{Résultats et livrables}
Modèle de données persistant et fonctionnel...

\subsection{Difficultés rencontrées et solutions}
...

\section{Synthèse de la phase}
Bilan des sprints 3-4, leçons apprises...

\section{Conclusion du chapitre}
...

% ============================================================================
% CHAPITRE 8: SPRINTS 5-6 - CRUD BACKEND ET FRONTEND DASHBOARD
% ============================================================================
\chapter{Sprints 5-6: CRUD Backend et Frontend Dashboard}
\label{chap:sprints5-6}

\section{Vue d'ensemble de la phase}
Objectifs globaux des semaines 9-12...

\section{Sprint 5: CRUD Backend Produits/Packs/Garanties (Semaines 9-10)}
\subsection{Objectifs du sprint}
Implémenter les opérations CRUD complètes...

\subsection{Planification}
\subsubsection{User Stories}
US-5.1: Implémenter CRUD Produit \\
US-5.2: Implémenter CRUD Pack \\
US-5.3: Implémenter CRUD Garantie \\
US-5.4: Implémenter association PackGarantie \\
US-5.5: Implémenter les endpoints REST \\
US-5.6: Ajouter validation et exceptions

\subsubsection{Estimation et priorisation}
...

\subsection{Réalisation}
\subsubsection{Service Produit}
CRUD complet, validation, logique métier...

\subsubsection{Service Pack}
Gestion des packs, association avec produits...

\subsubsection{Service Garantie}
CRUD garanties...

\subsubsection{Service PackGarantie}
Gestion des relations avec champs métier...

\subsubsection{Controllers REST}
Endpoints, documentation Swagger...

\subsection{Résultats et livrables}
API Backend CRUD complète...

\subsection{Difficultés rencontrées et solutions}
...

\section{Sprint 6: Frontend - Dashboard et Navigation (Semaines 11-12)}
\subsection{Objectifs du sprint}
Implémenter la structure Frontend, navigation, dashboard...

\subsection{Planification}
\subsubsection{User Stories}
US-6.1: Créer la structure Angular standalone \\
US-6.2: Implémenter la sidebar de navigation \\
US-6.3: Implémenter la navbar avec search \\
US-6.4: Créer le dashboard principal \\
US-6.5: Implémenter le breadcrumb dynamique \\
US-6.6: Configurer PrimeNG composants

\subsubsection{Estimation et priorisation}
...

\subsection{Réalisation}
\subsubsection{Structure Angular}
Modules, composants, routing...

\subsubsection{Sidebar de navigation}
Menu dynamique, icônes, responsive...

\subsubsection{Navbar avec search}
Barre de recherche, profil utilisateur, notifications...

\subsubsection{Dashboard principal}
Widgets, statistiques, vue d'ensemble...

\subsubsection{Breadcrumb dynamique}
Navigation hiérarchique...

\subsubsection{Intégration PrimeNG}
Composants UI, thèmes...

\subsection{Résultats et livrables}
Interface Frontend de base fonctionnelle...

\subsection{Difficultés rencontrées et solutions}
...

\section{Synthèse de la phase}
Bilan des sprints 5-6, leçons apprises...

\section{Conclusion du chapitre}
...

% ============================================================================
% CHAPITRE 9: SPRINTS 7-8 - FRONTEND GESTION ET IA
% ============================================================================
\chapter{Sprints 7-8: Frontend Gestion et Moteur IA}
\label{chap:sprints7-8}

\section{Vue d'ensemble de la phase}
Objectifs globaux des semaines 13-16...

\section{Sprint 7: Frontend - Gestion Produits/Packs/Garanties (Semaines 13-14)}
\subsection{Objectifs du sprint}
Implémenter les interfaces de gestion CRUD...

\subsection{Planification}
\subsubsection{User Stories}
US-7.1: Créer interface création/édition Produit \\
US-7.2: Créer interface création/édition Pack \\
US-7.3: Créer interface création/édition Garantie \\
US-7.4: Implémenter l'affichage hiérarchique Produit->Pack->Garantie \\
US-7.5: Ajouter les formulaires de validation \\
US-7.6: Implémenter l'affichage des champs PackGarantie

\subsubsection{Estimation et priorisation}
...

\subsection{Réalisation}
\subsubsection{Interface Produit}
Formulaire création/édition, liste, détails...

\subsubsection{Interface Pack}
Gestion des packs, association avec produits...

\subsubsection{Interface Garantie}
CRUD garanties...

\subsubsection{Affichage hiérarchique}
Vue Produit->Pack->Garantie avec champs PackGarantie...

\subsubsection{Formulaires de validation}
Validation reactive, messages d'erreur...

\subsection{Résultats et livrables}
Interfaces de gestion métier complètes...

\subsection{Difficultés rencontrées et solutions}
...

\section{Sprint 8: Moteur de Recommandation IA (Semaines 15-16)}
\subsection{Objectifs du sprint}
Implémenter le moteur de scoring et l'intégration IA...

\subsection{Planification}
\subsubsection{User Stories}
US-8.1: Intégrer Google Gemini API \\
US-8.2: Implémenter le moteur de scoring métier \\
US-8.3: Définir la formule de scoring \\
US-8.4: Implémenter l'analyse de profil client \\
US-8.5: Générer les prompts de recommandation \\
US-8.6: Implémenter l'extraction structurée des réponses \\
US-8.7: Retourner le Top 3 avec scores détaillés

\subsubsection{Estimation et priorisation}
...

\subsection{Réalisation}
\subsubsection{Intégration Google Gemini API}
Configuration client, gestion des erreurs...

\subsubsection{Moteur de scoring métier}
Implémentation de la formule de scoring:
\begin{equation}
score = (ageWeight \times 0.15) + (budgetWeight \times 0.20) + (beneficiaryWeight \times 0.15) + (medicalRiskWeight \times 0.20) + (guaranteeMatchWeight \times 0.20) + (coverageLevelWeight \times 0.10)
\end{equation}

\subsubsection{Analyse de profil client}
Extraction et analyse des informations client...

\subsubsection{Prompt engineering}
Structure des prompts pour Gemini...

\subsubsection{Extraction structurée}
Parsing des réponses IA en objets Java...

\subsubsection{Top 3 recommandations}
Génération des scores détaillés et explications...

\subsection{Résultats et livrables}
Moteur de recommandation IA fonctionnel...

\subsection{Difficultés rencontrées et solutions}
...

\section{Synthèse de la phase}
Bilan des sprints 7-8, leçons apprises...

\section{Conclusion du chapitre}
...

% ============================================================================
% CHAPITRE 10: SPRINTS 9-10 - CHATBOT ET DEVOPS
% ============================================================================
\chapter{Sprints 9-10: Chatbot Intelligente et DevOps}
\label{chap:sprints9-10}

\section{Vue d'ensemble de la phase}
Objectifs globaux des semaines 17-20...

\section{Sprint 9: Chatbot Intelligente et UI/UX (Semaines 17-18)}
\subsection{Objectifs du sprint}
Implémenter le chatbot et finaliser l'UI/UX...

\subsection{Planification}
\subsubsection{User Stories}
US-9.1: Créer l'interface chatbot type ChatGPT \\
US-9.2: Implémenter la différenciation user/assistant \\
US-9.3: Ajouter avatars et timestamps \\
US-9.4: Implémenter le typing indicator \\
US-9.5: Afficher les recommandations structurées \\
US-9.6: Implémenter le Light/Dark Mode \\
US-9.7: Ajouter la configuration conversationnelle

\subsubsection{Estimation et priorisation}
...

\subsection{Réalisation}
\subsubsection{Interface chatbot type ChatGPT}
Layout moderne, zone de messages, input...

\subsubsection{Différenciation user/assistant}
Styles distincts, avatars...

\subsubsection{Avatars et timestamps}
Affichage des métadonnées de message...

\subsubsection{Typing indicator}
Animation de chargement...

\subsubsection{Affichage des recommandations structurées}
Cartes pour Top 3, scores, explications...

\subsubsection{Light/Dark Mode}
Thème cohérent, bascule fonctionnelle...

\subsubsection{Configuration conversationnelle}
Mode création vs mode recommandation...

\subsection{Résultats et livrables}
Chatbot professionnel avec UI/UX moderne...

\subsection{Difficultés rencontrées et solutions}
...

\section{Sprint 10: DevOps, Tests et Finalisation (Semaines 19-20)}
\subsection{Objectifs du sprint}
Dockerisation, CI/CD, Monitoring, Tests...

\subsection{Planification}
\subsubsection{User Stories}
US-10.1: Dockeriser tous les services \\
US-10.2: Créer docker-compose.yml \\
US-10.3: Configurer GitHub Actions CI/CD \\
US-10.4: Implémenter tests JUnit backend \\
US-10.5: Implémenter tests Angular frontend \\
US-10.6: Configurer Prometheus monitoring \\
US-10.7: Créer dashboards Grafana \\
US-10.8: Ajouter documentation Swagger \\
US-10.9: Tests finaux et corrections bugs

\subsubsection{Estimation et priorisation}
...

\subsection{Réalisation}
\subsubsection{Dockerisation des services}
Dockerfiles pour Gateway, GestionProduit, Frontend...

\subsubsection{Docker Compose}
Orchestration de tous les services, réseau, volumes...

\subsubsection{CI/CD GitHub Actions}
Pipeline: build, test, Docker build, push...

\subsubsection{Tests Backend}
JUnit 5, Mockito, tests services et controllers...

\subsubsection{Tests Frontend}
Karma/Jasmine, tests composants critiques...

\subsubsection{Monitoring Prometheus}
Configuration Actuator, métriques...

\subsubsection{Dashboards Grafana}
CPU, RAM, temps de réponse, erreurs...

\subsubsection{Documentation Swagger}
Endpoints, modèles, exemples...

\subsubsection{Tests finaux}
Validation end-to-end, corrections bugs...

\subsection{Résultats et livrables}
Solution DevOps complète, tests, documentation...

\subsection{Difficultés rencontrées et solutions}
...

\section{Synthèse de la phase}
Bilan des sprints 9-10, leçons apprises...

\section{Conclusion du chapitre}
...

% ============================================================================
% CHAPITRE 11: RÉALISATION ET IMPLÉMENTATION
% ============================================================================
\chapter{Réalisation et Implémentation}
\label{chap:realisation}

\section{Environnement de développement}
\subsection{Outils et technologies}
IDE, JDK, Node.js, Docker, MongoDB...

\subsection{Configuration du projet}
Structure des projets Maven/NPM...

\section{Implémentation Backend}
\subsection{Service GestionProduit}
\subsubsection{Configuration Spring Boot}
application.yml, dépendances...

\subsubsection{Implémentation des entités}
Code Produit, Pack, Garantie, PackGarantie...

\subsubsection{Implémentation des repositories}
MongoRepository, requêtes personnalisées...

\subsubsection{Implémentation des services}
Logique métier, validation...

\subsubsection{Implémentation des controllers}
Endpoints REST, Swagger...

\subsection{Service GestionUser}
...

\subsection{API Gateway}
Configuration des routes, filtres...

\section{Implémentation Frontend}
\subsection{Structure Angular}
Modules, composants, services...

\subsection{Composants principaux}
\subsubsection{Dashboard}
Layout, widgets, navigation...

\subsubsection{Gestion Produits}
Liste, formulaire, détails...

\subsubsection{Gestion Packs}
Interface d'association...

\subsubsection{Chatbot}
Interface conversationnelle...

\subsection{Services Angular}
Communication avec backend, gestion d'état...

\subsection{Intégration PrimeNG}
Composants utilisés, customisation...

\subsection{Dark Mode}
Implémentation du thème...

\section{Implémentation de l'IA}
\subsection{Intégration Google Gemini API}
Configuration, client HTTP...

\subsection{Moteur de scoring}
Implémentation de l'algorithme...

\subsection{Prompt Engineering}
Structure des prompts...

\subsection{Extraction et parsing}
Traitement des réponses IA...

\section{Implémentation de la sécurité}
\subsection{Configuration Keycloak}
Realm, clients, rôles...

\subsection{Spring Security}
Configuration JWT, filtres...

\section{Implémentation DevOps}
\subsection{Dockerisation}
Dockerfiles, docker-compose...

\subsection{CI/CD GitHub Actions}
Workflow YAML...

\subsection{Monitoring}
Configuration Prometheus/Grafana...

\section{Tests}
\subsection{Tests Backend}
JUnit, Mockito, tests services...

\subsection{Tests Frontend}
Karma/Jasmine, tests composants...

\section{Conclusion du chapitre}
...

% ============================================================================
% CHAPITRE 12: RÉSULTATS ET DISCUSSION
% ============================================================================
\chapter{Résultats et Discussion}
\label{chap:resultats}

\section{Résultats fonctionnels}
\subsection{Gestion des produits}
Captures d'écran, démonstration...

\subsection{Gestion des packs et garanties}
Captures d'écran, hiérarchie affichée...

\subsection{Moteur de recommandation}
Exemples de recommandations, scores...

\subsection{Chatbot}
Captures d'écran, flux conversationnel...

\section{Résultats techniques}
\subsection{Performance}
Temps de réponse, charge...

\subsection{Sécurité}
Tests de sécurité, audit...

\subsection{Scalabilité}
Tests de charge...

\section{Validation des objectifs}
\subsection{Objectifs atteints}
Checklist des livrables...

\subsection{Écarts et justifications}
Ce qui n'a pas pu être fait et pourquoi...

\section{Discussion}
\subsection{Forces du projet}
Points forts de la solution...

\subsection{Faiblesses et limitations}
Contraintes, améliorations possibles...

\subsection{Comparaison avec l'existant}
Apport par rapport aux solutions existantes...

\section{Conclusion du chapitre}
...

% ============================================================================
% CHAPITRE 13: CONCLUSION ET PERSPECTIVES
% ============================================================================
\chapter{Conclusion et Perspectives}
\label{chap:conclusion}

\section{Synthèse du projet}
Résumé des réalisations, apprentissages...

\section{Apport personnel et professionnel}
Compétences acquises, expériences...

\section{Perspectives d'amélioration}
\subsection{Améliorations techniques}
RAG avancé, mémoire conversationnelle...

\subsection{Améliorations fonctionnelles}
Nouveaux cas d'utilisation, extensions...

\subsection{Améliorations DevOps}
Kubernetes, monitoring avancé...

\section{Conclusion générale}
Mot de fin...

% ============================================================================
% CHAPITRE 14: BIBLIOGRAPHIE
% ============================================================================
\chapter{Bibliographie}
\label{chap:bibliographie}

\section{Ouvrages}
...

\section{Articles scientifiques}
...

\section{Documentation technique}
...

\section{Ressources en ligne}
...

% ============================================================================
% ANNEXES
% ============================================================================
\appendix
\chapter{Annexes}
\label{chap:annexes}

\section{Annexe A: Code source}
\subsection{Extrait de code GestionProduit}
...

\subsection{Extrait de code Moteur de Scoring}
...

\section{Annexe B: Diagrammes complémentaires}
\subsection{Diagramme de séquence détaillé}
...

\subsection{Diagramme de déploiement}
...

\section{Annexe C: Configuration fichiers}
\subsection{application.yml}
...

\subsection{docker-compose.yml}
...

\subsection{GitHub Actions workflow}
...

\section{Annexe D: Manuel utilisateur}
Guide d'utilisation de l'application...

\section{Annexe E: Rapports de tests}
Résultats des tests unitaires et intégration...

\end{document}
